\documentclass[aspectratio=169]{beamer}

% ---------- ENCODING & FONTS ----------
\usepackage[utf8]{inputenc}
\usepackage[T1]{fontenc}

% ---------- THEME & LOOK ----------
\usetheme{Madrid}
\usecolortheme{seahorse}
\usefonttheme{professionalfonts}
\setbeamertemplate{navigation symbols}{}
\setbeamercovered{transparent}
\setbeamertemplate{itemize items}[circle]
\setbeamertemplate{enumerate items}[default]
\setlength{\parskip}{0.25em}

% ---------- COLORS ----------
% Beamer already loads xcolor, so we just define our colors
\definecolor{accent}{RGB}{0,92,184}
\definecolor{soft}{RGB}{233,244,255}
\definecolor{lightred}{RGB}{245,210,210}
\setbeamercolor{structure}{fg=accent}
\setbeamercolor{title}{fg=white,bg=accent}
\setbeamercolor{frametitle}{fg=accent}
\setbeamercolor{block title}{bg=accent!12,fg=accent}
\setbeamercolor{block body}{bg=soft}

% ---------- PACKAGES ----------
\usepackage{booktabs}
\usepackage{amsmath, amssymb}
\usepackage{array}
\usepackage{multirow}
\usepackage{hyperref}
\hypersetup{colorlinks=true,linkcolor=accent,urlcolor=accent}

% ---------- SAFE TRANSITION FALLBACK ----------
\makeatletter
\@ifundefined{transfade}{\newcommand{\transfade}[1][]{}}{}
\makeatother

% ---------- PLACEHOLDER BOX FOR FIGURES ----------
\newcommand{\placeholderbox}[2][0.82\linewidth]{%
  \begin{center}
    \fbox{\begin{minipage}[c][0.46\textheight][c]{#1}
      \centering \vspace{0.3em}
      \textbf{#2}\\[0.6em]
      \emph{Insert figure/diagram here later.}
    \end{minipage}}
  \end{center}
}

% ---------- TITLE ----------
\title{\textbf{Long-Run Economic Growth: Sources and Policies \vspace{.5cm}\\  }}
\subtitle{ {\large ECON 101 - Macroeconomics}}
%\institute{UNBC}
\author{Muhebullah Karimzada}
\date{Winter 2026}



\begin{document}

% ============================================================
% TITLE
% ============================================================
\begin{frame}[plain]
  \titlepage
\end{frame}

% ============================================================
\begin{frame}{Chapter Outline}
\transfade
\begin{itemize}[<+->]
  \item \textbf{7.1} The Importance of Economic Growth
  \item \textbf{7.2} The Economic Growth Model
  \item \textbf{7.3} New Growth Theory and Knowledge Capital
  \item \textbf{7.4} Why Isn\'t the Whole World Rich?
\end{itemize}
\end{frame}

% ============================================================
\section{7.1 The Importance of Economic Growth}
% ============================================================

\begin{frame}{Obtaining Economic Growth}
\transfade
\begin{itemize}[<+->]
  \item Previous chapter: how to \textbf{measure} growth (real GDP, real GDP per capita).
  \item This chapter: how to \textbf{explain} and \textbf{influence} long-run growth.
  \item Key idea:
  \begin{itemize}[<+->]
    \item Economic growth is \alert{not inevitable}.
    \item History includes long periods of stagnation with little or no increase in output per person.
  \end{itemize}
  \item Central questions:
  \begin{itemize}[<+->]
    \item Why have some countries achieved rapid, sustained growth?
    \item Why have others failed to keep up and remained poor?
  \end{itemize}
\end{itemize}
\end{frame}

\begin{frame}{The Industrial Revolution}
\transfade
\begin{itemize}[<+->]
  \item \textbf{Industrial Revolution:}
  \begin{itemize}[<+->]
    \item Application of mechanical power to the production of goods and services.
    \item Began in England around 1750.
  \end{itemize}
  \item Before that:
  \begin{itemize}[<+->]
    \item Production relied mainly on human and animal muscle power.
    \item Very slow (almost zero) growth in real GDP per capita.
  \end{itemize}
  \item Mechanical power (steam engines, later electricity) allowed:
  \begin{itemize}[<+->]
    \item England and later the United States, France, Germany, and others
    \item to experience sustained \textbf{long-run economic growth}.
  \end{itemize}
\end{itemize}
\end{frame}

\begin{frame}{Table 7.1 -- The Effects of Different Growth Rates on Living Standards (Placeholder)}
\transfade
\placeholderbox{Table 7.1 -- The Effects of Different Growth Rates on Living Standards.\\[0.4em]
\textbf{Instruction:} Reproduce the exact table from the textbook here, with the same headings and numerical values showing how different constant growth rates (e.g., 1\%, 2\%, 3\%) change income levels over time (e.g., 35, 70 years).}
\end{frame}

\begin{frame}{Why Do Growth Rates Matter?}
\transfade
\begin{itemize}[<+->]
  \item Small differences in annual growth rates lead to \alert{huge} differences in living standards over decades.
  \item A country that grows too slowly fails to raise living standards:
  \begin{itemize}[<+->]
    \item Not just about smartphones and TVs,
    \item but about basic indicators:
    \begin{itemize}[<+->]
      \item infant mortality,
      \item life expectancy,
      \item access to clean water, health care, and education.
    \end{itemize}
  \end{itemize}
  \item Example contrast:
  \begin{itemize}[<+->]
    \item In many high-income countries, fewer than 6 of 1{,}000 babies die before age one.
    \item In some of the poorest countries, more than 50 of 1{,}000 die before age one.
  \end{itemize}
\end{itemize}
\end{frame}

\begin{frame}{Differences in Incomes Across Countries}
\transfade
\begin{itemize}[<+->]
  \item Economists often divide economies into:
  \begin{itemize}[<+->]
    \item \textbf{High-income (industrial) countries:}
    \begin{itemize}[<+->]
      \item e.g., Australia, Canada, Japan, New Zealand, the United States, Western Europe.
    \end{itemize}
    \item \textbf{Developing (low- and middle-income) countries:}
    \begin{itemize}[<+->]
      \item many countries in Africa, Asia, and Latin America.
    \end{itemize}
  \end{itemize}
  \item Some economies, like Singapore, South Korea, and Taiwan, have become \textbf{newly industrializing countries} with very high growth rates.
  \item Real GDP per capita differs dramatically even after adjusting for cost-of-living differences.
  \item Example from 2020:
  \begin{itemize}[<+->]
    \item very high income: Monaco (over \$170{,}000 per person),
    \item very low income: Burundi (under \$300 per person).
  \end{itemize}
\end{itemize}
\end{frame}

\begin{frame}{Figure 7.1 -- Real GDP per Capita Around the World (Diagram Placeholder)}
\transfade
\placeholderbox{Figure 7.1 -- Map or bar chart of real GDP per capita across countries, showing large gaps between high-income, newly industrializing, and developing countries.}
\end{frame}

% ============================================================
\section{7.2 The Economic Growth Model}
% ============================================================

\begin{frame}{7.2 Learning Objective}
\transfade
\begin{block}{Goal}
Use the economic growth model to explain why growth rates differ across countries.
\end{block}
\begin{itemize}[<+->]
  \item An \textbf{economic growth model} explains growth rates in \textbf{real GDP per capita} over the long run.
  \item Key idea from Chapter 6:
  \begin{itemize}[<+->]
    \item Long-run growth in real GDP per person is driven by growth in \textbf{labour productivity}.
  \end{itemize}
\end{itemize}
\end{frame}

\begin{frame}{Labour Productivity and Its Determinants}
\transfade
\begin{itemize}[<+->]
  \item \textbf{Labour productivity:}
  \[
    \text{Labour productivity} =
    \frac{\text{Quantity of goods and services produced}}{\text{Number of workers or hours worked}}.
  \]
  \item Two main factors determine labour productivity:
  \begin{enumerate}[<+->]
    \item \textbf{Quantity of capital available to workers}
    \item \textbf{Level of technology}
  \end{enumerate}
  \item Our model focuses on:
  \begin{itemize}[<+->]
    \item changes in the quantity of capital per worker, and
    \item \textbf{technological change}: a change in the amount of output a firm can produce with a given quantity of inputs.
  \end{itemize}
\end{itemize}
\end{frame}

\begin{frame}{Sources of Technological Change}
\transfade
\begin{itemize}[<+->]
  \item Three main sources of technological change:
  \begin{enumerate}[<+->]
    \item \textbf{Better machinery and equipment}
    \item \textbf{Increases in human capital}
    \item \textbf{Better organization and management of production}
  \end{enumerate}
\end{itemize}
\end{frame}

\begin{frame}{Better Machinery and Equipment}
\transfade
\begin{itemize}[<+->]
  \item Innovations such as:
  \begin{itemize}[<+->]
    \item the steam engine,
    \item machine tools,
    \item electric generators,
    \item computers and robots,
    \item advanced telecommunications.
  \end{itemize}
  \item These allow each worker to produce \alert{more} per hour with the same effort.
  \item Example:
  \begin{itemize}[<+->]
    \item A clerk with a modern computer system processes many more transactions than one using only paper records.
  \end{itemize}
\end{itemize}
\end{frame}

\begin{frame}{Increases in Human Capital}
\transfade
\begin{itemize}[<+->]
  \item \textbf{Human capital:} accumulated knowledge and skills workers acquire from:
  \begin{itemize}[<+->]
    \item education,
    \item training,
    \item and experience.
  \end{itemize}
  \item More human capital means:
  \begin{itemize}[<+->]
    \item workers can adopt new technologies more quickly,
    \item they can solve problems more effectively,
    \item they are more flexible when the economy changes.
  \end{itemize}
\end{itemize}
\end{frame}

\begin{frame}{Better Organization and Management}
\transfade
\begin{itemize}[<+->]
  \item Sometimes productivity rises not because of new machines, but because of:
  \begin{itemize}[<+->]
    \item better management,
    \item improved organization of production,
    \item more efficient supply chains.
  \end{itemize}
  \item Example:
  \begin{itemize}[<+->]
    \item Just-in-time inventory systems reduce waste and allow firms to respond quickly to demand, raising output per worker.
  \end{itemize}
\end{itemize}
\end{frame}

% ------------------------------------------------------------
% Per-Worker Production Function
% ------------------------------------------------------------

\begin{frame}{The Per-Worker Production Function}
\transfade
\begin{itemize}[<+->]
  \item To formalize these ideas, we use a \textbf{per-worker production function}:
  \begin{itemize}[<+->]
    \item It shows the relationship between \textbf{real GDP per hour worked} and \textbf{capital per hour worked}, holding technology constant.
  \end{itemize}
  \item With a fixed level of technology:
  \begin{itemize}[<+->]
    \item The first units of capital are the most effective.
    \item Additional units increase output, but by \alert{smaller and smaller} amounts.
  \end{itemize}
  \item This is called \textbf{diminishing returns} to capital.
\end{itemize}
\end{frame}

\begin{frame}{Figure 7.3 -- The Per-Worker Production Function (Diagram Placeholder)}
\transfade
\placeholderbox{Figure 7.3 -- Per-worker production function: real GDP per hour worked vs. capital per hour worked, with diminishing returns.}
\end{frame}

\begin{frame}{Diminishing Returns to Capital}
\transfade
\begin{itemize}[<+->]
  \item \textbf{Diminishing returns:}
  \begin{itemize}[<+->]
    \item Increasing one input (capital) progressively, while holding other inputs constant, yields smaller and smaller increases in output.
  \end{itemize}
  \item Intuition:
  \begin{itemize}[<+->]
    \item Give a worker their first machine or computer: huge productivity gain.
    \item Give them the tenth very similar machine: smaller extra gain.
  \end{itemize}
  \item This means that increasing capital per worker alone cannot generate high growth forever.
\end{itemize}
\end{frame}

% ------------------------------------------------------------
% Technological Change and the Production Function
% ------------------------------------------------------------

\begin{frame}{Technological Change and the Production Function}
\transfade
\begin{itemize}[<+->]
  \item When technology improves, the entire per-worker production function shifts \textbf{upward}.
  \item This means:
  \begin{itemize}[<+->]
    \item For any given amount of capital per worker,
    \item output per worker is higher.
  \end{itemize}
  \item Technological change is therefore crucial for sustained long-run growth.
\end{itemize}
\end{frame}

\begin{frame}{Figure 7.4 -- Technological Change Increases Output (1 of 2) (Diagram Placeholder)}
\transfade
\placeholderbox{Figure 7.4 (1 of 2) -- For a country with low capital per worker, increases in capital are very effective at raising real GDP per hour worked.}
\end{frame}

\begin{frame}{Figure 7.4 -- Technological Change Increases Output (2 of 2) (Diagram Placeholder)}
\transfade
\placeholderbox{Figure 7.4 (2 of 2) -- With high capital per worker, further increases in capital face strong diminishing returns; technological change becomes the main source of growth.}
\end{frame}

\begin{frame}{Capital Deepening vs. Technological Change}
\transfade
\begin{itemize}[<+->]
  \item \textbf{Capital deepening:}
  \begin{itemize}[<+->]
    \item An increase in the amount of capital per worker.
  \end{itemize}
  \item In relatively poor countries:
  \begin{itemize}[<+->]
    \item Capital deepening can produce large gains in productivity.
  \end{itemize}
  \item In relatively rich countries:
  \begin{itemize}[<+->]
    \item Capital deepening still helps, but diminishing returns set in.
    \item Continuing growth in living standards requires \textbf{sustained technological change}.
  \end{itemize}
\end{itemize}
\end{frame}

% ============================================================
\section{7.3 New Growth Theory and Knowledge Capital}
% ============================================================

\begin{frame}{New Growth Theory}
\transfade
\begin{itemize}[<+->]
  \item The growth model we have used so far is associated with \textbf{Robert Solow} (1950s).
  \item In the Solow model:
  \begin{itemize}[<+->]
    \item Productivity growth is key to long-run growth.
    \item But technological change is treated as something that happens \textbf{outside} the model (an exogenous factor).
  \end{itemize}
  \item \textbf{New growth theory}, developed by \textbf{Paul Romer} and others:
  \begin{itemize}[<+->]
    \item Emphasizes that technological change is driven by \textbf{economic incentives}.
    \item It is an outcome of decisions made by firms, researchers, and governments within the market system.
  \end{itemize}
\end{itemize}
\end{frame}

\begin{frame}{Knowledge Capital and Economic Growth}
\transfade
\begin{itemize}[<+->]
  \item Romer argues that accumulation of \textbf{knowledge capital} is a key determinant of growth.
  \item \textbf{Knowledge capital:}
  \begin{itemize}[<+->]
    \item The stock of ideas, scientific knowledge, designs, and blueprints.
    \item Generated by research and development (R\&D) and other innovative activities.
  \end{itemize}
  \item Increases in knowledge capital arise from:
  \begin{itemize}[<+->]
    \item R\&D in firms and universities,
    \item learning-by-doing,
    \item spillovers from one firm or country to another.
  \end{itemize}
\end{itemize}
\end{frame}

\begin{frame}{Physical Capital vs. Knowledge Capital}
\transfade
\begin{itemize}[<+->]
  \item \textbf{Physical capital:}
  \begin{itemize}[<+->]
    \item Rival and excludable (a private good).
    \item One machine can only be used by one worker at a time.
    \item Leads to \textbf{diminishing returns}.
  \end{itemize}
  \item \textbf{Knowledge capital:}
  \begin{itemize}[<+->]
    \item Nonrival and often nonexcludable (a public good).
    \item One firm\'s use of an idea does not prevent others from using it.
    \item At the level of the whole economy, this can generate \textbf{increasing returns} to scale.
  \end{itemize}
\end{itemize}
\end{frame}

\begin{frame}{Free-Rider Problem and Underinvestment in Knowledge}
\transfade
\begin{itemize}[<+->]
  \item Because knowledge capital is a public good, firms may not capture the full benefits of their R\&D.
  \item \textbf{Free-rider problem:}
  \begin{itemize}[<+->]
    \item People or firms can benefit from an idea without paying for it.
  \end{itemize}
  \item Example:
  \begin{itemize}[<+->]
    \item A major lab develops a transistor or a new drug.
    \item Other firms learn from this and earn profits, even though they did not bear the original research cost.
  \end{itemize}
  \item Result:
  \begin{itemize}[<+->]
    \item Market economy produces \textbf{too little} knowledge capital relative to the socially optimal amount.
  \end{itemize}
\end{itemize}
\end{frame}

\begin{frame}{Government\'s Role in Supporting Knowledge Capital}
\transfade
\begin{itemize}[<+->]
  \item Governments can help increase the accumulation of knowledge capital in three ways:
  \begin{enumerate}[<+->]
    \item Protecting intellectual property
    \item Supporting research and development
    \item Subsidizing education
  \end{enumerate}
\end{itemize}
\end{frame}

\begin{frame}{Protecting Intellectual Property}
\transfade
\begin{itemize}[<+->]
  \item \textbf{Patents:}
  \begin{itemize}[<+->]
    \item Exclusive right to produce a product for a fixed period (e.g., 20 years from the filing date).
    \item Allow firms to earn profits from their inventions, increasing their incentive to innovate.
  \end{itemize}
  \item \textbf{Copyrights:}
  \begin{itemize}[<+->]
    \item Similar protection for creative works (books, films, music).
    \item Typically last for the creator\'s life plus many years after.
  \end{itemize}
  \item Policy trade-off:
  \begin{itemize}[<+->]
    \item Give inventors enough protection to encourage R\&D,
    \item while ensuring that society eventually benefits from widespread use of the ideas.
  \end{itemize}
\end{itemize}
\end{frame}

\begin{frame}{Supporting R\&D and Subsidizing Education}
\transfade
\begin{itemize}[<+->]
  \item \textbf{Supporting R\&D:}
  \begin{itemize}[<+->]
    \item Governments conduct research in public labs.
    \item They subsidize R\&D at universities and private firms through grants and tax incentives.
  \end{itemize}
  \item \textbf{Subsidizing education:}
  \begin{itemize}[<+->]
    \item High-quality education helps create a workforce capable of doing research and adopting new technologies.
    \item Public funding of schooling and universities reduces barriers to human capital accumulation.
  \end{itemize}
\end{itemize}
\end{frame}

% ============================================================
\section{7.4 Why Isn\'t the Whole World Rich?}
% ============================================================

\begin{frame}{7.4 Learning Objective}
\transfade
\begin{block}{Goal}
Explain economic catch-up and why many poor countries have not experienced rapid growth.
\end{block}
\begin{itemize}[<+->]
  \item The economic growth model predicts that \textbf{poor countries should grow faster} than rich ones.
  \item This would imply that poor countries eventually \textbf{catch up} to the income levels of rich countries.
\end{itemize}
\end{frame}

\begin{frame}{Catch-Up: The Convergence Hypothesis}
\transfade
\begin{itemize}[<+->]
  \item \textbf{Catch-up (convergence):}
  \begin{itemize}[<+->]
    \item Prediction that GDP per capita in poor countries will grow faster than in rich countries.
  \end{itemize}
  \item Why might we expect this?
  \begin{itemize}[<+->]
    \item Additional capital is more productive in countries with small capital stocks.
    \item Poor countries can adopt technologies already used in rich countries, rather than inventing everything themselves.
  \end{itemize}
\end{itemize}
\end{frame}

\begin{frame}{Figure 7.6 -- Catch-Up Predicted by the Model (Diagram Placeholder)}
\transfade
\placeholderbox{Figure 7.6 -- Hypothetical graph showing real GDP per capita in poor and rich countries converging over time as poor countries grow faster.}
\end{frame}

\begin{frame}{Evidence of Catch-Up Among High-Income Countries}
\transfade
\begin{itemize}[<+->]
  \item When we look only at high-income countries, we see evidence of catch-up:
  \begin{itemize}[<+->]
    \item Countries that were relatively poor in 1960 (e.g., Hong Kong, South Korea) had faster growth.
    \item Countries that were relatively rich in 1960 (e.g., the United States, Switzerland) had slower growth.
  \end{itemize}
  \item So among high-income countries, convergence appears to have occurred over the postwar period.
\end{itemize}
\end{frame}

\begin{frame}{Figure 7.7 -- Catch-Up Among High-Income Countries (Diagram Placeholder)}
\transfade
\placeholderbox{Figure 7.7 -- Growth from 1960 to recent years plotted against initial real GDP per capita for high-income countries, showing a negative relationship (poorer grew faster).}
\end{frame}

\begin{frame}{Most of the World Has Not Been Catching Up}
\transfade
\begin{itemize}[<+->]
  \item When we extend the analysis to \textbf{all} countries, the simple catch-up prediction performs poorly.
  \item Many low-income countries:
  \begin{itemize}[<+->]
    \item have grown much more slowly than high-income countries,
    \item and some have even experienced periods of negative growth.
  \end{itemize}
  \item So we need to understand:
  \begin{itemize}[<+->]
    \item Why some poor countries grow quickly and others do not.
  \end{itemize}
\end{itemize}
\end{frame}

\begin{frame}{Figure 7.8 -- Most of the World Has Not Been Catching Up (Diagram Placeholder)}
\transfade
\placeholderbox{Figure 7.8 -- Scatter plot of growth vs. initial income for many countries, showing weak or no convergence once low-income countries are included.}
\end{frame}

\begin{frame}{Why Have High-Income Countries Stopped Catching Up to the U.S.?}
\transfade
\begin{itemize}[<+->]
  \item In 1990, many high-income countries had real GDP per capita somewhat below the U.S. level.
  \item By 2019, many had \textbf{not} closed the gap and some had fallen further behind.
  \item Possible reasons:
  \begin{itemize}[<+->]
    \item U.S. labour markets are relatively \textbf{flexible}.
    \item U.S. workers tend to:
    \begin{itemize}[<+->]
      \item enter the labour force earlier,
      \item retire later.
    \end{itemize}
    \item The U.S. financial system is relatively \textbf{efficient} and \textbf{liquid}, making it attractive for investment.
    \item Access to venture capital and support for start-ups is strong, helping new firms grow.
  \end{itemize}
\end{itemize}
\end{frame}

\begin{frame}{Figure 7.9 -- Other High-Income Countries vs. the U.S. (Diagram Placeholder)}
\transfade
\placeholderbox{Figure 7.9 -- Bars showing real GDP per capita relative to the U.S. in 1990 and in 2019 for selected high-income countries; most red bars are equal or lower than blue bars.}
\end{frame}

\begin{frame}{Why Are Many Poor Countries Growing So Slowly?}
\transfade
\begin{itemize}[<+->]
  \item Economists often highlight four key factors:
  \begin{enumerate}[<+->]
    \item Failure to enforce the rule of law
    \item Wars and revolutions
    \item Poor public education and health
    \item Low rates of saving and investment
  \end{enumerate}
\end{itemize}
\end{frame}

\begin{frame}{Failure to Enforce the Rule of Law}
\transfade
\begin{itemize}[<+->]
  \item \textbf{Rule of law:}
  \begin{itemize}[<+->]
    \item Government\'s ability to enforce the laws, particularly with respect to:
    \begin{itemize}[<+->]
      \item protecting \textbf{private property},
      \item enforcing \textbf{contracts}.
    \end{itemize}
  \end{itemize}
  \item \textbf{Property rights:}
  \begin{itemize}[<+->]
    \item Individuals and firms must have the exclusive right to use their property and to buy or sell it.
  \end{itemize}
  \item Without strong rule of law:
  \begin{itemize}[<+->]
    \item entrepreneurs fear expropriation and corruption,
    \item they invest less in new businesses and technologies.
  \end{itemize}
\end{itemize}
\end{frame}

\begin{frame}{Wars and Revolutions}
\transfade
\begin{itemize}[<+->]
  \item Wars and internal conflicts:
  \begin{itemize}[<+->]
    \item destroy physical capital,
    \item disrupt education and health systems,
    \item create uncertainty that discourages investment.
  \end{itemize}
  \item Example:
  \begin{itemize}[<+->]
    \item A country experiencing civil war may have the same real GDP per capita decades apart.
    \item Once peace is restored, growth can resume at high rates.
  \end{itemize}
\end{itemize}
\end{frame}

\begin{frame}{Poor Public Education and Health}
\transfade
\begin{itemize}[<+->]
  \item Weak public schools and limited access to education:
  \begin{itemize}[<+->]
    \item reduce \textbf{human capital},
    \item make it harder for workers to adopt new technologies.
  \end{itemize}
  \item Poor health care and inadequate nutrition:
  \begin{itemize}[<+->]
    \item lower worker productivity,
    \item reduce the number of hours people are able to work.
  \end{itemize}
  \item Investing in education and health is both:
  \begin{itemize}[<+->]
    \item a social policy,
    \item and a long-run growth policy.
  \end{itemize}
\end{itemize}
\end{frame}

\begin{frame}{Low Rates of Saving and Investment}
\transfade
\begin{itemize}[<+->]
  \item Many low-income countries have:
  \begin{itemize}[<+->]
    \item underdeveloped and insecure financial systems,
    \item limited access to credit for households and firms.
  \end{itemize}
  \item This creates a \textbf{vicious cycle}:
  \begin{itemize}[<+->]
    \item Low income \(\rightarrow\) low saving \(\rightarrow\) low investment \(\rightarrow\) slow growth \(\rightarrow\) low income.
  \end{itemize}
  \item Breaking this cycle often requires:
  \begin{itemize}[<+->]
    \item reforms to financial institutions,
    \item support from international organizations and foreign investors,
    \item domestic policies that encourage saving and investment.
  \end{itemize}
\end{itemize}
\end{frame}

% ============================================================
\section*{Chapter Summary}
% ============================================================

\begin{frame}{Summary}
\transfade
\begin{itemize}[<+->]
  \item Long-run economic growth is crucial for improving living standards.
  \item Growth models emphasize:
  \begin{itemize}[<+->]
    \item labour productivity,
    \item capital per worker,
    \item and technological change.
  \end{itemize}
  \item New growth theory places \textbf{knowledge capital} and \textbf{incentives} at the centre of growth.
  \item Governments can promote growth by:
  \begin{itemize}[<+->]
    \item protecting intellectual property,
    \item supporting R\&D,
    \item and subsidizing education.
  \end{itemize}
  \item While some countries have caught up to high-income levels, many have not, due to:
  \begin{itemize}[<+->]
    \item weak rule of law,
    \item conflict,
    \item poor education and health,
    \item and low saving and investment.
  \end{itemize}
\end{itemize}
\end{frame}

\begin{frame}
\centering
\Huge \textbf{Thank you!}
\end{frame}

\end{document}
